% Dirichlet Round 07 English TeX
% Translation from G. Lejeune Dirichlet, Werke II, pp. 89--96.
% Source span checked against Werke II scan images 102--109.
\documentclass[11pt]{article}
\usepackage[a4paper,margin=1in]{geometry}
\usepackage[T1]{fontenc}
\usepackage[utf8]{inputenc}
\usepackage{lmodern}
\usepackage{amsmath,amssymb}
\usepackage{microtype}
\setlength{\parindent}{1.2em}
\setlength{\parskip}{0.25em}
\newcommand{\leg}[2]{\left(\frac{#1}{#2}\right)}
\begin{document}

\begin{center}
{\Large G. Lejeune Dirichlet's Works. II.}\par\vspace{1.2em}
{\Large VIII.}\par\vspace{0.5em}
{\LARGE \textbf{On the Representability of Numbers}}\par
{\LARGE \textbf{by Three Squares.}}\par\vspace{1em}
{\large By G. Lejeune Dirichlet.}\par\vspace{0.5em}
{\small Crelle, Journal für die reine und angewandte Mathematik, vol. 40, pp. 228--232.}
\end{center}

\begin{center}
{\large \textbf{On the Representability of Numbers by Three Squares.}}
\end{center}

The theory of decomposing the integers $n$ which are not of one of the forms $4k$, $8k+7$ into three squares without a common divisor belongs to the more complicated parts of higher arithmetic, if one wishes to develop this theory completely, that is, not merely to prove representability, but also to determine the number of all possible decompositions. This number can be expressed either by means of the binary forms of determinant $-n$, as Gauss first showed\footnote{Disquisitiones arithmeticae, art. 229.}, or also independently of the latter\footnote{Crelle's Journal, vol. 21, p. 155. The editor of the Werke adds: p. 496 of vol. I of this edition of G. Lejeune Dirichlet's works. K.}. Since, however, there are cases in which only representability is to be assumed, as for example in the proof given by Cauchy\footnote{Exercices de Mathématiques par Cauchy, année 1826, p. 265. The editor of the Werke adds: Oeuvres complètes d'Augustin Cauchy, IIe Série, Tome VI, p. 320. K.}, later simplified by Legendre, of Fermat's theorem on polygonal numbers, which rests solely on the fact that every number, apart from those excluded above, is the sum of three squares, a simple proof of representability seems not to be without interest. Such a proof is to be communicated in this paper.

For this purpose one first needs the known theorem that every positive ternary form of determinant $-1$, that is, every expression
\begin{equation}\tag{1}
 ax^2+by^2+cz^2+2a'yz+2b'xz+2c'xy,
\end{equation}
whose integral coefficients satisfy the equation
\begin{equation}\tag{2}
 aa'^2+bb'^2+cc'^2-abc-2a'b'c'=-1
\end{equation}
and, moreover, satisfy the conditions that
\[
 a,\quad b,\quad c,\quad bc-a'^2,\quad ac-b'^2,\quad ab-c'^2
\]
are positive (of which conditions, incidentally, the first and fourth imply the others), is equivalent to the form
\begin{equation}\tag{3}
 x^2+y^2+z^2.
\end{equation}
To prove this theorem it suffices to verify that, for determinant $-1$, the expression (3) is the only reduced positive form.

If the form (1) is to be reduced, then, by the theorem proved at the end of the preceding paper\footnote{Vol. II, p. 48 of this edition of G. Lejeune Dirichlet's works. The cited paper, ``Über die Reduction der positiven quadratischen Formen mit drei unbestimmten ganzen Zahlen'', is printed in vol. 40 of Crelle's Journal immediately before this paper. K.}, $abc$ must not be greater than $2$. Hence, if one assumes $a\leq b\leq c$, it follows at once that
\[
 a=b=1.
\]
On the other hand, by the definition of reduced forms, both $2c'$ and $2b'$, apart from sign, must not be greater than $a$, and $2a'$ must not be greater than $b$; thus one obtains
\[
 a'=b'=c'=0,
\]
and then, from (2),
\[
 c=1.
\]

This being assumed, the decomposability of the positive number $a$ will have been shown as soon as one has found a positive ternary form (1) of determinant $-1$ whose first coefficient is $a$. For since such a form is equivalent to the form (3), the latter can be transformed into it, and one obtains
\[
 a=\alpha^2+\alpha'^2+\alpha''^2,
\]
where $\alpha$, $\alpha'$, $\alpha''$ are three of the nine substitution coefficients and cannot have a common divisor, since every term of the determinant formed from these nine coefficients has either $\alpha$, or $\alpha'$, or $\alpha''$ as a factor, and this determinant is equal to unity.

Thus everything comes down to the following: regarding $a$ as given, satisfy equation (2) by five integers $b$, $c$, $a'$, $b'$, $c'$, with only the single remaining condition that $bc-a'^2$ must become positive. If one takes $b'=1$, $c'=0$, the equation becomes
\[
 b=a\Delta-1,
\]
where
\[
 \Delta=bc-a'^2
\]
is set; and it remains only to show that the positive number $\Delta$ can be chosen so that $-\Delta$ is a quadratic residue of $b=a\Delta-1$, since under this assumption $c$ and $a'$ can be so determined that $a'^2-bc=-\Delta$.

We shall now show that the condition just stated can always be satisfied by an odd value of $\Delta$ such that, if $a$ has the form $4k+2$, then $b=a\Delta-1$ is an odd prime number $p$, and if $a$ is odd but not of the form $8k+7$, then $b$ is twice such an odd prime number $p$.

We begin with the second case. We have the equation
\[
 2p=a\Delta-1,
\]
where for the moment $p$ is not yet assumed prime, but only odd. Put $\Delta=8t+\varepsilon$, where $\varepsilon$ is one of the numbers $1$, $3$, $5$, $7$. One sees at once that in each of the four cases presented by $a$ with respect to the divisor $8$, $\Delta$ admits two forms with respect to this divisor, or $\varepsilon$ admits two values, if $p$, as required, is to be odd. For each of the eight cases thus obtained, apply the law of reciprocity to the left side of the equation following from $2p=a\Delta-1$,
\[
 \leg{p}{\Delta}=\leg{-2}{\Delta},
\]
where the Legendre symbol is used in the extended sense introduced by Jacobi; put for $\leg{-2}{\Delta}$ its known value, and then multiply by $\leg{-1}{p}=\pm1$, where $\pm1$ will likewise be known according to the corresponding linear form of $p$. One obtains the following results:
\[
\begin{array}{r@{\quad}l}
 a=8k+1, &
 \begin{cases}
 \Delta=8t+3,\quad p=4s+1,\quad \leg{-\Delta}{p}=+1,\\
 \Delta=8t+7,\quad p=4s+3,\quad \leg{-\Delta}{p}=-1,
 \end{cases}\\[1.3em]
 a=8k+3, &
 \begin{cases}
 \Delta=8t+1,\quad p=4s+1,\quad \leg{-\Delta}{p}=+1,\\
 \Delta=8t+5,\quad p=4s+3,\quad \leg{-\Delta}{p}=+1,
 \end{cases}\\[1.3em]
 a=8k+5, &
 \begin{cases}
 \Delta=8t+3,\quad p=4s+3,\quad \leg{-\Delta}{p}=+1,\\
 \Delta=8t+7,\quad p=4s+1,\quad \leg{-\Delta}{p}=-1,
 \end{cases}\\[1.3em]
 a=8k+7, &
 \begin{cases}
 \Delta=8t+1,\quad p=4s+3,\quad \leg{-\Delta}{p}=-1,\\
 \Delta=8t+5,\quad p=4s+1,\quad \leg{-\Delta}{p}=-1.
 \end{cases}
\end{array}
\]

It follows that, if, as assumed, $a$ is not of the form $8k+7$, the condition $\leg{-\Delta}{p}=1$ can always be met. That in all eight cases $p$ can also be made a prime number is clear from the expression
\[
 p=\frac12(a\Delta-1)=4at+\frac12(a\varepsilon-1),
\]
in which $\frac12(a\varepsilon-1)$ is, by the choice of $\varepsilon$, odd and also has no common divisor with $a$. This expression is therefore the general term of an arithmetic progression which necessarily contains prime numbers. Once one has a prime number $p$ for which $\leg{-\Delta}{p}=1$, $-\Delta$ is a quadratic residue of $p$ and consequently also of $2p$, as was required.

In the case of an even $a$ we at once put $a=4k+2$, since one knows beforehand that for $a=4k$ the condition cannot be satisfied. Since in the equation
\[
 p=a\Delta-1
\]
we assume $\Delta$ to be odd, $p$ has the form $4s+1$, and we obtain
\[
 \leg{-1}{\Delta}=\leg{p}{\Delta}=\leg{\Delta}{p}=\leg{-\Delta}{p}.
\]
Thus $\Delta$ must be given the form $4t+1$ in order that $\leg{-\Delta}{p}=1$. One obtains
\[
 p=4at+a-1,
\]
and this expression can again become a prime number, for which $-\Delta$ is then a quadratic residue.

The application of the procedure just developed is not restricted to the case of determinant $-1$, and we add a second example for determinant $-3$.

If one again seeks, as above, the reduced forms for this determinant, then the condition
\[
 abc\leq 6
\]
gives, under the assumption $a\leq b\leq c$, the value $1$ for $a$, while $b$ can be equal to $1$ or to $2$. It follows, since $2c'$ and $2b'$ must not be numerically greater than $a=1$, that
\[
 b'=c'=0,
\]
and the equation by which the determinant is determined reduces to
\[
 a'^2-bc=-3.
\]
Now, since $2a'$ also must not be numerically greater than $b$, if $b=1$ then $a'=0$, and consequently $c=3$. If, on the other hand, $b=2$, then $2a'=0$ or $2a'=\pm2$. The first value does not satisfy the above equation, in which $bc=4$. It remains only to set $a'=\pm1$, whence $c=2$. Neglecting the lower sign, which amounts only to changing the sign of $z$, one therefore obtains the two reduced forms
\begin{equation}\tag{4}
 x^2+y^2+3z^2,\qquad x^2+2y^2+2z^2+2yz,
\end{equation}
so that every positive ternary form (1) of determinant $-3$, that is, one in which the coefficients satisfy the equation
\begin{equation}\tag{5}
 aa'^2+bb'^2+cc'^2-abc-2a'b'c'=-3,
\end{equation}
will be equivalent to one of these forms (4). If one now examines the residues modulo $8$ which the expressions (4) leave when the elements $x$, $y$, $z$ are given all combinations of even and odd values, excluding only the case of three even values simultaneously, since the elements are always to be assumed without a common divisor, one finds that the first of the forms (4) offers all possible residues modulo $8$, whereas the second expression assumes neither of the forms $8k+5$, $4k$ for any combination. If one further considers that two equivalent forms always represent the same numbers, then, for a form of determinant $-3$ which represents a number $8k+5$ or $4k$ (which is in particular always the case when one of its first three coefficients, for instance $a$, is such a number), one can at once conclude that it is not equivalent to the second form, and consequently that it is equivalent to the first of the forms (4).

It remains to prove that the given number $a$ is representable by the first of the forms (4). By the preceding method, and by again putting in (5)
\[
 b'=1,\qquad c'=0,\qquad \Delta=bc-a'^2,
\]
we need only show that the positive number $\Delta$ can be so determined that $b=a\Delta-3$ has the form $8k+5$ or $4k$, and that $-\Delta$ is at the same time a quadratic residue of $b$. We restrict ourselves here to the numbers $a$ which have no common divisor with the determinant, that is, which are not divisible by $3$. Put
\[
 \Delta=8\Delta',
\]
where $\Delta'$ is to be odd and not divisible by $3$. Then $b$ will be relatively prime to $\Delta'$ and will have the form $8k+5$, and only the other condition remains to be fulfilled. From the equation $b=8a\Delta'-3$ it follows at once that
\[
 \leg{b}{\Delta'}=\leg{-3}{\Delta'},
\]
or, if one takes account of the congruence $b\equiv1\pmod 4$ and applies known theorems,
\[
 \leg{b}{\Delta'}=\leg{\Delta'}{b}=\leg{-\Delta'}{b}=\leg{-3}{\Delta'}.
\]
Multiplication by $\leg{8}{b}=\leg{2}{b}=-1$ gives from this
\[
 \leg{-\Delta}{b}=-\leg{-3}{\Delta'},
\]
and one sees that the condition $\leg{-\Delta}{b}=1$ is fulfilled if
\[
 \Delta'=6t-1
\]
and consequently
\[
 b=48at-8a-3
\]
is set. Since this latter expression can evidently be made equal to a prime number, it is proved that every number not divisible by $3$ can be expressed by the form
\[
 x^2+y^2+3z^2
\]
in such a way that $x$, $y$, $z$ have no common divisor. Similar considerations can be applied to the numbers divisible by $3$, for whose representability an easily obtained condition is required, as well as to the numbers expressible by the second of the forms (4).

\begin{flushright}
Berlin, June 1850.
\end{flushright}

\end{document}
